% SPDX-License-Identifier: LPPL-1.3c
% Copyright (C) 2026 Christophe Jorssen
% Author and Current Maintainer: Christophe Jorssen <christophe.jorssen@gmail.com>
% LPPL maintenance status: maintained
% This file is part of luafemm. See LICENSE and LICENSES.md for its terms.
\section{A rotor moves inside a circular air band}
\label{sec:sliding}

\subsection{Independent mechanical positions}
A rotor-position study holds the stator excitation fixed and compares
fields, torque and flux linkage at independent mechanical angles. Redrawing and remeshing the entire machine at every
angle would mix physical changes with changes in the finite-element mesh.
A \emph{moving air band} keeps the bulk triangles attached to their respective
bodies. Only a narrow annulus of air receives new triangles at each position.

The band treats independent \emph{planar static or fixed-angle harmonic} positions.
The field is solved afresh at each angle, with nonlinear laws where supported. It introduces
neither angular velocity nor motion-induced voltage. With the default |method=p1|, its triangles are ordinary
P1 finite elements, not FEMM's specialised analytical air-gap elements (AGE).
There can still be discretisation ripple when the band diagonals change.
The alternative analytical method and study commands are described in
Section~\ref{sec:machine-studies}.

\Needspace{10\baselineskip}
\subsection{Declare once, capture, then choose positions}
The order matters. First declare the problem, materials, circuits and all
physical paths or nodes. Then capture a mesh-only template. Finally create
one new problem per requested angle and perform the measurements there.
\begin{enumerate}
\item Draw every rotor domain strictly inside an inner circle. Draw every
  stator domain outside an outer circle. Leave a source-free air collar on
  both sides of the circles: neither trace may coincide with a material edge.
\item Attach |femm/sliding band| to the declaration picture. The key captures
  at the end of that picture, after all its paths. Do not request a field,
  mesh, profile value, measurement or operating point before this capture.
\item In a new picture, set |femm/problem={sliding band=machine,rotor angle=30}|.
  This copies the physical data, turns the rotor and connects the band.
\item Draw the inherited geometry/field and measure torque, circuit linkage
  or ordinary profiles. Repeat step 3 for another angle.
\end{enumerate}

\begin{codeexample}[code only]
\tikzset{machine position/.style={
  femm/problem={sliding band=machine,rotor angle=#1},
  femm/cache={file=machine-angle-#1,mode=auto}}}
% A preceding declaration picture captured the template named machine.
\begin{tikzpicture}[machine position=30]
  \pic {femm geometry};
  \pic {femm field};
  \femmmeasure[name=at-30,quantity=force,method=band,parameter=30]
  \femmmeasure[name=port-30,quantity=circuit,circuit=drive,parameter=30]
  \path[femm/profile={name=gap-30,samples=181}]
    (4.8,0) arc[start angle=0,end angle=360,radius=4.8];
\end{tikzpicture}
\end{codeexample}
The picture transformation controls presentation. It does not turn the physical
rotor: that is the job of |rotor angle|. In particular, an untagged rotated
ellipse is just a drawing overlay. Tagged physical paths in the declaration
picture retain their full TikZ expressiveness; the usual flattening tolerance
still controls their computational polygons. The band traces themselves are
regular polygonal approximations of circles.

\Needspace{10\baselineskip}
\subsection{A complete reluctance-machine experiment}
This example uses an elliptical rotor and an annular stator with two prescribed
100-turn winding legs. The signs of their turns close the winding, so its
linkage is independent of the potential gauge. The steel table is illustrative
and single-valued. The two field panels use independently chosen contour levels;
compare the measured linkage, rather than visible line counts, for magnitude.
The dashed circles show the artificial band, not a new material interface.

The source starts with a relatively coarse mesh. While adjusting colours and
labels, use only two or three angles and |mesh size=1.2|. Then decrease
the bulk spacing, increase |segments| and check several |layers| values.
For a final quantitative plot, also move the exterior box farther away.
A smooth-looking interpolated plot is not evidence of converged torque.

\luafemmexamplebegin{sliding-band}

Use the following styles and geometry for the position study. The capture and measurement loops produce data; the later pictures display those data.

% manual example: sliding-band/1

\begin{codeexample}[setup code]
\definecolor{fluxblue}{RGB}{20,85,135}
\definecolor{torquered}{RGB}{185,65,45}
\newbox\MachineCalculation
\def\Mesh{.8}
\tikzset{
  femm/materials/core/.style={bh={0/0,.25/25.293,.5/59.375,
    .75/146.191,1/400,1.25/1040.527,1.5/2428.125,
    1.75/5098.926,2/9800,2.5/29546.875,3/73200}},
  winding/.style={femm/region={material=copper,circuit=drive,
    turns=#1,winding fill=.7,mesh size=.4}},
  position/.style={femm/problem={sliding band=machine,rotor angle=#1},
    femm/cache={file=\jobname-angle-#1,mode=auto}},
  femm/arrow spacing=7mm}
\setbox\MachineCalculation=\hbox{
  \begin{tikzpicture}[femm/problem={xmin=-12,xmax=12,ymin=-12,ymax=12,
      mesh size=\Mesh,depth=10,interpolation=femm},
    femm/circuit={name=drive,current=10},
    femm/cache={file=\jobname-template,mode=auto},
    femm/sliding band={name=machine,inner radius=4.3,outer radius=5.3,
      segments=96,layers=3}]
    \path[femm/region={material=core,mesh size=.4}]
      (0,0) ellipse[x radius=3.8,y radius=1.8];
    \path[even odd rule,femm/region={material=core,mesh size=.4}]
      (0,0) circle[radius=9] (0,0) circle[radius=6];
    \path[winding=100] (-8.3,-1.2) rectangle (-7.2,1.2);
    \path[winding=-100] (7.2,-1.2) rectangle (8.3,1.2);
  \end{tikzpicture}}
\end{codeexample}

The following definitions complete the shared setup.

% manual example: sliding-band/2

\begin{codeexample}[setup code]
\foreach \Angle in {0,15,...,180}{
  \setbox\MachineCalculation=\hbox{
    \begin{tikzpicture}[position=\Angle]
      \femmmeasure[name=torque-\Angle,quantity=force,method=band,parameter=\Angle]
      \femmmeasure[name=port-\Angle,quantity=circuit,circuit=drive,parameter=\Angle]
    \end{tikzpicture}}}
\end{codeexample}

The following picture uses the preceding definitions.

% manual example: sliding-band/3

\begin{codeexample}[width=0pt]
\foreach \Angle in {30,75}{
  \begin{tikzpicture}[position=\Angle,scale=2.3]
    \clip (-11,-10.2) rectangle (11,11.5);
    \pic {femm geometry};
    \begin{scope}
      \clip (-10,-10) rectangle (10,9.5);
      \pic[femm/lines=23,femm/field style={fluxblue}] {femm field};
    \end{scope}
    \draw[teal!65!black,densely dashed,thin]
      (0,0) circle[radius=4.3] (0,0) circle[radius=5.3];
    \node[fill=white,font=\small] at (0,10.3) {$\theta=\Angle^\circ$};
  \end{tikzpicture}}
\end{codeexample}

% manual example: sliding-band/4

\begin{codeexample}[setup code]
\def\TorqueResults{torque-0,torque-15,torque-30,torque-45,torque-60,torque-75,
  torque-90,torque-105,torque-120,torque-135,torque-150,torque-165,torque-180}
\def\PortResults{port-0,port-15,port-30,port-45,port-60,port-75,port-90,
  port-105,port-120,port-135,port-150,port-165,port-180}
\end{codeexample}

The following picture uses the preceding definitions.

% manual example: sliding-band/5

\begin{codeexample}[width=0pt]
\begin{tikzpicture}
\begin{axis}[width=7cm,height=5.8cm,grid=major,xmin=0,xmax=180,
  xtick={0,45,90,135,180},xlabel={Mechanical angle (degrees)},
  ylabel={Torque (mN\,m)},y filter/.expression={1000*y},scaled y ticks=false,
  title={Alignment torque}]
  \femmresultplot[results={\TorqueResults},component=torque,
    plot options={torquered,thick,mark=*,mark options={fill=torquered}}]
\end{axis}
\end{tikzpicture}
\end{codeexample}

% manual example: sliding-band/6

\begin{codeexample}[width=0pt]
\begin{tikzpicture}
\begin{axis}[width=7cm,height=5.8cm,grid=major,xmin=0,xmax=180,
  xtick={0,45,90,135,180},xlabel={Mechanical angle (degrees)},
  ylabel={Flux linkage (mWb)},y filter/.expression={1000*y},scaled y ticks=false,
  title={The magnetic path changes}]
  \femmresultplot[results={\PortResults},component=linkage,
    plot options={fluxblue,thick,mark=square*}]
\end{axis}
\end{tikzpicture}
\end{codeexample}

\luafemmexampleend{sliding-band}

At fixed current, the rotor tends towards increasing magnetic coenergy. For
this example it aligns approximately at $90^\circ$ and the torque reverses
there. The geometry repeats after $180^\circ$; the API itself normalises angles
only modulo $360^\circ$. Asymmetry in a discretised template can cause small
differences between nominally equivalent half-turn positions.

The loop in this example is ordinary |\foreach|. Each named result remains
available to |\femmresultplot| after its picture has ended. Store the intended
abscissa explicitly with |parameter|. Changing |femm/lines|, colour, scale or
profile sample count does not change the solved physical position.

\Needspace{10\baselineskip}
\subsection{Capture keys and geometric contract}
\Needspace{7\baselineskip}
\begin{key}{/tikz/femm/sliding band=\marg{key list}}
Schedule capture at the end of this picture, binding it to this picture's model.
The keys belong to |/femm/sliding band|. There is no default template name.
Only one band is supported per template. All physical declarations must precede
the capture and no mesh or solve may already exist. A mesh cache restored
\emph{by the capture itself} is allowed.
\end{key}
\Needspace{7\baselineskip}
\begin{command}{\femmslidingband\oarg{key list}}
Capture immediately. Place this command after the last physical declaration.
It accepts the same key family as the picture key. It does not draw or solve.
Each invocation resets the capture options to the defaults below before applying
its list; package styles can supply recurring settings.
\end{command}
\Needspace{7\baselineskip}
\begin{key}{/femm/sliding band/name=\meta{name} (initially empty)}
Required document-wide template identifier. Duplicate names are errors. The
name remains available across TeX groups for the current run; it is not itself
a saved file or a persistent object identifier.
\end{key}
\Needspace{7\baselineskip}
\begin{key}{/femm/sliding band/model=\meta{name} (initially empty)}
For the immediate command, select an existing named model instead of the current
one. It must still be unmeshed. The picture-level key supplies its own originating
model name after the user's list, so later pictures cannot redirect its capture.
\end{key}
\Needspace{7\baselineskip}
\begin{key}{/femm/sliding band/inner radius=\meta{number} (initially 0)}
Inner trace radius in model units. Required positive value. Every rotating path
must lie strictly inside the inscribed polygon, not merely inside the circle.
\end{key}
\Needspace{7\baselineskip}
\begin{key}{/femm/sliding band/outer radius=\meta{number} (initially 0)}
Outer trace radius in model units, strictly larger than the inner radius.
Every stator path must lie outside this circle; compound stator paths may
surround it but their material must not fill the band. The entire circle must
lie strictly inside the rectangular problem box.
\end{key}
\Needspace{7\baselineskip}
\begin{key}{/femm/sliding band/center x=\meta{number} (initially 0)}
Horizontal coordinate of the rotation centre in model units.
\end{key}
\Needspace{7\baselineskip}
\begin{key}{/femm/sliding band/center y=\meta{number} (initially 0)}
Vertical coordinate of the same centre. These are model coordinates, unaffected
by later drawing transformations. They do not set the torque origin.
\end{key}
\Needspace{7\baselineskip}
\begin{key}{/femm/sliding band/segments=\meta{integer} (initially 96)}
Number of vertices on each initially declared trace and each intermediate ring;
allowed range 16--1024. The template mesher may subdivide the trace chords,
so their final vertex counts can differ. Raising this value refines the angular
resolution and the polygonal geometry; it does not refine the bulk interiors.
\end{key}
\Needspace{7\baselineskip}
\begin{key}{/femm/sliding band/layers=\meta{integer} (initially 2)}
Number of radial strips, from 1 to 16. Intermediate rings are equally spaced
in radius and fixed in the stator frame. With inner radius $r_i$, outer radius
$r_o$, $L$ layers and $N$ segments, the nesting check requires
\[
 \bigl(r_i+(r_o-r_i)/L\bigr)\cos(\pi/N)>r_i.
\]
If it fails, increase |segments| or reduce |layers|. Excessive radial refinement
with insufficient angular refinement creates poor triangles. The ordinary
minimum-angle and maximum-area targets, when nonzero, are checked on every
new position; an unmet target is an error, not a silently accepted mesh.
\end{key}
All six numerical capture keys accept ordinary PGF mathematical expressions.
No Lua source is evaluated from key values. A path crossing or filling the band
is rejected, including a large air refinement rectangle. Use local spacing
on rotor/stator regions instead. Both traces must be in air with vacuum
permeability, zero coercivity and no current. The region order retains its
usual meaning inside each bulk domain.

\Needspace{10\baselineskip}
\subsection{Position and measurement keys}
\Needspace{7\baselineskip}
\begin{key}{/femm/problem/sliding band=\meta{name} (initially empty)}
Create an independent position from the named template. An empty value selects
the ordinary problem workflow. The template owns geometry, materials, circuits,
exterior conditions, mesh targets, units, depth and outer-solver settings. Do not
redeclare them in a position picture. Linear-solver settings are inherited but
may be explicitly overridden using Section~\ref{sec:linear-controls}. Graphics, profiles, selections and
measurements may still be added. Cache options belong to the new picture.
\end{key}
\Needspace{7\baselineskip}
\begin{key}{/femm/problem/rotor angle=\meta{number} (initially 0)}
Mechanical rotation in degrees, counterclockwise about the band centre, relative
to the captured geometry. PGF expressions are allowed, for example |90/3|.
Negative values are allowed and finite angles are reduced modulo $360^\circ$.
A nonzero angle without |sliding band| is an error. It is not an electrical
angle or a display phase. Ordinary TeX grouping controls problem-key inheritance.
\end{key}
The existing |/femm/problem/circuit currents| list can override prescribed
series currents in amperes. Unlisted currents retain their template values.
This differs from the zero-excitation default of an incremental response.
For a harmonic template, |circuit phases={drive=30}| also overrides the temporal
phase in degrees; unlisted phases retain their template values. Frequency is
owned by the template. Unknown circuit names are rejected. Current, phase and
mechanical angle enter the physical cache signature. No other physical problem override
is accepted, even if it would numerically equal the inherited setting.

\Needspace{7\baselineskip}
\begin{key}{/femm/measure/method=band}
Additional choice for |quantity=force| on a sliding position. Integrate Maxwell
stress across the entire air band to obtain the total |Fx|, |Fy| and |torque|
on everything inside it. No body selection or profile is needed or accepted.
The origin remains |origin x=0,origin y=0| in model units; explicitly set it to
the band centre for torque about an offset rotation axis. Forces use newtons,
torque uses newton metres, and positive torque is counterclockwise. The problem
must have positive |depth|. Existing integral components cover the whole model.
\end{key}
The general measurement default remains |method=weighted|. Set |method=band|
explicitly, or define a style:
\begin{codeexample}[code only]
\pgfkeys{/femm/measure/rotor torque/.style={quantity=force,method=band}}
\femmmeasure[rotor torque,name=at-60,parameter=60]
\end{codeexample}
|weighting| and |tolerance| configure the ordinary Laplace-weighted body method;
they do not change the band's prescribed weights. For an independent diagnostic,
register a closed contour strictly inside source-free air and use the existing
|method=maxwell,profile=...|. A single contour can have more radial bias than
the volume integral. Compare several contours and refined meshes.

Ordinary profiles and circuit measurements retain their syntax. For example,
|\femmplot[profile=gap-30,component=By,abscissa=angle]| inside a PGFPlots axis
plots the angular gap field. Set the profile centre explicitly for an offset
machine. Material selections use the original material names even when a
rotated anisotropic rotor internally needs its own tensor copy. Such a material
filter also selects stator regions with that material; add a geometric selection
when isolating a body for the ordinary weighted-force method.

\Needspace{10\baselineskip}
\subsection{Numerical construction and accuracy}
The template is meshed once with two protected polygonal traces. Band triangles
are discarded; retained bulk nodes and triangle connectivity are copied into
each position. Rotor coordinates, material axes and permanent-magnet directions
rotate physically. Uniform, radial, tangential and angular magnetisation laws
follow their material frame, including the law centre. Stator coordinates stay
unchanged. Each new annular strip is triangulated by connecting its ordered
inner and outer vertices. Orientation predicates, edge incidence, coverage,
Euler characteristic and retained constraints are checked before solving.

For |method=band|, the nodal weight $w$ is one on the rotor bulk, zero on the
stator bulk, and $1-\ell/L$ on intermediate ring $\ell$. P1 interpolation makes
its gradient constant in each band triangle. With planar Maxwell stress
\[
 \mathbf T=\frac1{\mu_0}\left(\mathbf B\mathbf B^{\mathsf T}
             -\frac12|\mathbf B|^2\mathbf I\right),\qquad
 \mathbf F=-\int_{\rm band}\mathbf T\nabla w\,\mathrm dV,
\]
the torque is the moment of the same force density about the chosen origin.
The piecewise-constant field and linear moment arm are integrated exactly for
the discrete P1 field. This exact integration does not eliminate finite-element
error. The band's geometric weights avoid an additional Laplace solve and
average stress over its thickness. They are not FEMM's AGE torque formula.

Useful checks include preservation of a uniform field for arbitrary angles,
zero net stress in that uniform field, a magnet in an imposed transverse field,
and the fixed-current identity
$T=\partial W'/\partial\theta$ with $\theta$ in radians. Compare the energy
finite difference at several angular steps; making the step extremely small
can amplify reconnection ripple. Refine bulk spacing, angular segmentation,
radial layers and exterior distance independently. The implementation does not
promise noise-free cogging torque or an error bound from a small solver residual.

\Needspace{10\baselineskip}
\subsection{Persistence, interchange and current limits}
Use separate cache files for the template mesh and each angle/current state.
The template can reuse a mesh-only cache; each position can reuse its solution.
Frozen mode checks the geometry, sources and exact inherited mesh just as it
does for other models. Native format 24 includes a moving-band signature carrying
method identity, template provenance, normalised angle, retained/reconnected
mesh and force weights. Incompatible cache versions require recomputation. Template names must be recreated in each TeX run; they are not
serialised as a document-wide object registry.

A supported position exports as an ordinary static or supported linear harmonic |.fem| or |.ans| snapshot
at its current angle. Its material paths, magnet directions and mesh are already
rotated. A FEMM user can solve that one |.fem| normally; the file does not contain
a live moving-band controller or native AGE records. The answer export preserves
the computed P1 mesh and potentials. Existing representability checks still
apply: a rotated off-diagonal material tensor cannot be encoded in FEMM's
global diagonal material records, and nonlinear ANS export needs
|interpolation=femm|. Angular coercive laws retain their documented export rules.
Importing a snapshot creates an ordinary model, not a sliding template.

The scope is one complete circular band in a rectangular native planar
static or harmonic problem with |mesher=delaunay| and |field model=open|. Custom domains,
tagged/paired boundaries, periodic sectors, imported templates, callbacks,
parallel/voltage circuits and ideal circuits are excluded. A solved DC position
may supply an operating point; its responses retain that exact angle and mesh.
A harmonic template cannot contain permanent-magnet DC sources: use a separate
DC bias and incremental response instead. A native machine node may supply rotor/stator domains
when all its paths satisfy the air-collar contract. Existing TikZ node anchors
are drawing objects and are not recreated or rebound by a later position.

\Needspace{10\baselineskip}
\subsection{Lua entry points and validation}
\begin{codeexample}[code only]
local S = require("luafemm-sliding")
local template = S.capture(model, {
  inner_radius=4.3, outer_radius=5.3,
  center_x=0, center_y=0, segments=96, layers=3
})
local position = S.new(template, {
  angle=30, currents={drive=10},
  cache={file="machine-30",mode="auto"}
})
require("luafemm").solve(position)
local force = require("luafemm-post").measure(position, {
  name="at-30",quantity="force",method="band"
})
\end{codeexample}
|capture| meshes the supplied unmeshed model and returns independently owned
template data; treat those returned data as read-only. |new| returns an independent
meshed, unsolved model. Its angle defaults to zero, currents to an empty override
list and cache mode to |off|. Unknown options are errors. Physical edits to a
position after creation are rejected; create another position from the template.

The portable tests cover bulk identity, angle wrapping, affine fields, analytic
magnet torque, nonlinear coenergy differences, tensor and angular-law rotation,
source overrides and passive caches. The format matrix exercises shared cold,
warm and frozen data in plain LuaTeX, LuaLaTeX and ConTeXt MkIV. Independent validation exports static angle snapshots and compares torque
and linkage with independent xfemm remeshes and solves. Numerical evidence and
its tolerances are recorded in \nolinkurl{doc/validation.md}.


\subsection{AC torque and the first variation about a bias}
At a fixed rotor angle of $30^\circ$, distinguish two excitation models. An AC
current with no bias produces a mean torque and a twice-frequency ripple.
A small AC current added to a DC operating point produces a first-order torque
at the excitation frequency. A temporal phase of $30^\circ$ is independent of
the mechanical angle, even when the numbers happen to be equal.

The following example declares two templates: one at zero frequency for the DC
bias, and one at $50\,\mathrm{Hz}$ for the independent AC calculation. Each owns
its frequency. Solve the DC position and capture its operating point
\emph{after} turning the rotor. Create the small harmonic response from
that bias, without specifying another rotor angle or band. To study another
angle, first construct and solve another DC position. A response cannot
be rotated away from its bias mesh.

The material is deliberately linear and lossless so that the two waveforms
can be checked analytically: at equal peak current, the AC mean torque is half
the DC torque, while a small increment gives
$\delta T=2T_0\delta I/I_0$. A real nonlinear material instead needs its actual
differential tensor. Start with a few angles and a coarse mesh; refine the bulk,
angular traces and radial layers independently before interpreting small torques.
The two field panels use separate contour levels and do not imply equal flux.

\luafemmexamplebegin{machine-response}

% manual example: machine-response/1

\begin{codeexample}[setup code]
\newbox\ResponseCalculation
\tikzset{femm/materials/machine iron/.style={mu r=300},
  machine winding/.style={femm/region={material=copper,circuit=drive,turns=#1}},
  response field/.style={femm/lines=19,femm/arrows=false,
    femm/field style={blue!65!black}}}
\pgfkeys{/femm/measure/rotor load/.style={quantity=force,method=band}}
\foreach \Frequency in {0,50}{
  \setbox\ResponseCalculation=\hbox{
    \begin{tikzpicture}[femm/problem={xmin=-10,xmax=10,ymin=-10,ymax=10,
        mesh size=1,depth=10,frequency=\Frequency},
      femm/circuit={name=drive,current=10},
      femm/cache={file=\jobname-template-\Frequency,mode=auto},
      femm/sliding band={name=machine-\Frequency,
        inner radius=3.5,outer radius=4.5,segments=64,layers=2}]
      \path[femm/region={material=machine iron,mesh size=.5}]
        (0,0) ellipse[x radius=3,y radius=1.5];
      \path[machine winding=100] (-7,-2) rectangle (-6,2);
      \path[machine winding=-100] (6,-2) rectangle (7,2);
    \end{tikzpicture}}}
\end{codeexample}

The following picture uses the preceding definitions.

% manual example: machine-response/2

\begin{codeexample}[width=0pt]
\begin{tikzpicture}[femm/problem={sliding band=machine-0,rotor angle=30},
    femm/cache={file=\jobname-bias,mode=auto},
    femm/operating point={name=machine-bias},scale=2.5]
  \pic {femm geometry};
  \pic[response field] {femm field};
  \femmmeasure[rotor load,name=dc-load]
  \draw[teal,densely dashed] (0,0) circle[radius=4];
  \node[font=\small,fill=white] at (0,9) {DC bias: $I_0=10\,\mathrm A$};
\end{tikzpicture}
\end{codeexample}

% manual example: machine-response/3

\begin{codeexample}[width=0pt]
\begin{tikzpicture}[femm/problem={operating point=machine-bias,frequency=50,
    circuit currents={drive=1},circuit phases={drive=30}},
    femm/cache={file=\jobname-response,mode=auto},scale=2.5]
  \pic {femm geometry};
  \pic[response field,femm/phase=0] {femm field};
  \femmmeasure[rotor load,name=small-load]
  \node[font=\small,fill=white] at (0,9) {$\widehat{\delta I}=1e^{i30^\circ}\,\mathrm A$};
\end{tikzpicture}
\end{codeexample}

% manual example: machine-response/4

\begin{codeexample}[setup code]
\setbox\ResponseCalculation=\hbox{
  \begin{tikzpicture}[femm/problem={sliding band=machine-50,rotor angle=30,
      circuit phases={drive=30}},femm/cache={file=\jobname-ac,mode=auto}]
    \femmmeasure[rotor load,name=ac-load]
  \end{tikzpicture}}
\end{codeexample}

The following picture uses the preceding definitions.

% manual example: machine-response/5

\begin{codeexample}[width=0pt]
\begin{tikzpicture}
\begin{axis}[width=6.3cm,height=4.8cm,font=\small,xmin=0,xmax=360,xtick={0,90,180,270,360},
    grid=major,xlabel={Excitation phase (degrees)},ylabel={$T$ (mN\,m)},
    y filter/.expression={1000*y},scaled y ticks=false,title={AC without bias}]
  \femmforceplot[result=ac-load,component=torque,
    plot options={blue!65!black,thick,mark=none}]
\end{axis}
\end{tikzpicture}
\end{codeexample}

% manual example: machine-response/6

\begin{codeexample}[width=0pt]
\begin{tikzpicture}
\begin{axis}[width=6.3cm,height=4.8cm,font=\small,xmin=0,xmax=360,xtick={0,90,180,270,360},
    grid=major,xlabel={Excitation phase (degrees)},ylabel={$\delta T$ (mN\,m)},
    y filter/.expression={1000*y},scaled y ticks=false,title={Response around DC bias}]
  \femmforceplot[result=small-load,component=torque,
    plot options={orange!80!black,thick,mark=none}]
\end{axis}
\end{tikzpicture}
\end{codeexample}

\luafemmexampleend{machine-response}

The code composes a normal |/femm/measure| style for rotor loads. No new command
or plotting engine is required. |circuit phases| rotates excitation phasors;
|femm/phase| selects a temporal field view; |rotor angle| rotates geometry,
material axes and attached sources. With conducting rotor material these are
\emph{fixed-position} diffusion solves. No velocity, slip frequency, motional
$\mathbf v\times\mathbf B$ term, or back-emf from continuous rotation is implied.

\subsection{An eccentric rotor}
\label{sec:eccentric}
A static eccentricity moves the rotor centre away from the stator centre.
Choose a full-circle |method=p1| template. At each position the rotor first
rotates about the reference centre, then translates in the fixed stator frame.
Both its material axes and its magnetisation follow the rotation. Intermediate
air rings interpolate between the translated inner centre and the fixed outer
centre. All bulk triangles retain their connectivity.

\begin{key}{/femm/problem/rotor offset x=\meta{number} (initially 0)}
Horizontal rotor translation, in model units, after |rotor angle| is applied.
Requires a full-circle P1 sliding template. It affects the solve, geometry,
profiles and cache identity. A nonzero offset with a concentric Fourier/native
AGE method, or without a sliding template, is an error.
\end{key}
\begin{key}{/femm/problem/rotor offset y=\meta{number} (initially 0)}
Vertical translation with the same scope and units. For static eccentricity,
keep the two offsets fixed while changing angle. For an offset rotating with
the shaft, set them to expressions such as |e*cos(Angle)| and |e*sin(Angle)|
inside a TikZ loop. These are independent stationary solutions, not a transient
motion calculation.
\end{key}

\luafemmexamplebegin{eccentric-gap}
\begin{codeexample}[width=0pt]
\tikzset{femm/materials/eccentric iron/.style={mu r=100}}
\begin{tikzpicture}[
  femm/problem={xmin=-8,xmax=8,ymin=-8,ymax=8,mesh size=1,depth=10},
  femm/sliding band={name=eccentric-ring,inner radius=3,
    outer radius=4,segments=64,layers=2},
  scale=2.5,baseline=(current bounding box.center)]
  \draw[femm/region={material=eccentric iron}]
    (0,0) ellipse[x radius=2.6,y radius=1.2];
  \draw[femm/region={material=copper,current density=1e7}]
    (-6,-1) rectangle (-5,1);
  \draw[femm/region={material=copper,current density=-1e7}]
    (5,-1) rectangle (6,1);
\end{tikzpicture}
\begin{tikzpicture}[
  femm/problem={sliding band=eccentric-ring,rotor angle=25,
    rotor offset x=.45,rotor offset y=.2},
  scale=2.5,baseline=(current bounding box.center)]
  \pic {femm geometry};
  \pic[femm/lines=13] {femm field};
  \draw[teal,dashed] (.45,.2) circle[radius=3];
  \draw[teal,dashed] (0,0) circle[radius=4];
  \femmmeasure[name=eccentric-force,quantity=force,method=band,
    origin x=.45,origin y=.2]
\end{tikzpicture}
\end{codeexample}
\luafemmexampleend{eccentric-gap}

The dashed circles show how the band narrows on one side. Torque here is about
the displaced rotor centre; omitting the origin keys would instead measure
about $(0,0)$. Force uses newtons and torque uses newton metres. Existing profile
paths and |\femmresultplot| can display their variation over a position sweep.

The offset magnitude must be strictly smaller than the radial clearance.
A stronger polygonal nesting test applies to every layer. Every assembled
triangle must have positive orientation, the complete mesh must cover the
original outer domain, and all bulk constraints must survive. Requested minimum
angles and maximum areas are checked again after reconnecting the band.
A rejected mesh is not solved or cached. Increase angular segmentation, reduce
layers, or reduce eccentricity when polygonal layers overlap; simply relaxing
a solver tolerance cannot repair this geometry.

Start with a coarse bulk mesh, then refine it and the trace segmentation.
Check force against the displacement derivative of coenergy at fixed currents,
and torque against the angular derivative about the displaced centre. Repeat
with a smaller finite-difference step: changes of band diagonals can otherwise
obscure convergence. Offset values and the complete position mesh are part of
the cache signature.


For a time-dependent field with material history, use the distinct state and step
interface in Section~\ref{sec:time-motion}. Its full-circle P1 contract retains
conducting bulk nodes and supports prescribed motion or rotational dynamics.
